\begin{tikzpicture}[
    >=Stealth,
    every node/.style={font=\scriptsize},
    seed/.style={
        circle, draw=blue!70!black, fill=blue!18,
        line width=0.7pt, text width=2.0cm, align=center,
        inner sep=2pt, minimum size=20mm,
    },
    ent/.style={
        circle, draw=teal!70!black, fill=teal!12,
        line width=0.5pt, text width=1.6cm, align=center,
        inner sep=2pt, minimum size=14mm,
    },
    text2/.style={
        rectangle, rounded corners=4pt,
        draw=orange!75!black, fill=orange!10,
        line width=0.4pt, text width=2.4cm, align=center,
        inner sep=3pt, font=\tiny,
    },
    edge/.style={->, draw=black!55, line width=0.4pt},
    legendbox/.style={
        rectangle, draw=black!30, fill=white,
        inner sep=4pt,
    },
]
% 种子节点
\node[seed] (seed) at (0,0) {\textbf{TextNode}\\\scriptsize NERC LL20231101\\\scriptsize 电流差动保护逻辑};

% 1-hop 实体节点（围绕种子径向放置）
\node[ent] (e1) at ( 70:3.4) {电流差动\\保护};
\node[ent] (e2) at ( 25:3.4) {NERC\\LL20231101};
\node[ent] (e3) at (-25:3.4) {fault};
\node[ent] (e4) at (-70:3.4) {trip};
\node[ent] (e5) at (-130:3.4) {间歇性\\退化信号};
\node[ent] (e6) at ( 130:3.4) {230\,kV\\transmission};

% 2-hop 文本节点
\node[text2] (t1) at ($(e1.north) + (0,1.2)$) {方向比较闭锁\\保护中…};
\node[text2] (t2) at ($(e3.east) + (1.4,0)$) {故障清除后\\电机重新加速…};
\node[text2] (t3) at ($(e6.west) + (-1.4,0)$) {LL20231101\\手动替代值定义…};

% 边：seed → 1-hop entities
\foreach \e in {e1,e2,e3,e4,e5,e6} {
    \draw[edge] (seed) -- (\e);
}
% 边：1-hop entity → 2-hop text
\draw[edge] (e1) -- (t1);
\draw[edge] (e3) -- (t2);
\draw[edge] (e6) -- (t3);

% 图例（散布节点 + 简单标签）
\node[legendbox, anchor=north west, font=\tiny, text width=5.6cm,
      align=left] at (-6.4,-4.2)
{%
\tikz[baseline=-0.6ex]\node[seed, minimum size=4mm, text width=0pt, inner sep=0pt]{};\ 种子文本节点（top-$k$ 检索） \\
\tikz[baseline=-0.6ex]\node[ent,  minimum size=4mm, text width=0pt, inner sep=0pt]{};\ 1 跳实体节点（BFS 扩展） \\
\tikz[baseline=-0.6ex]\node[text2, text width=4mm,  inner sep=1pt]{};\ \ 2 跳文本节点（BFS 扩展） \\
\textcolor{black!55}{$\rightarrow$}\ REFERENCES 关系（实体引用源文档）%
};
\end{tikzpicture}
